\documentclass[11pt,a4paper]{article}
\usepackage[T1]{fontenc}
\usepackage[utf8]{inputenc}
\usepackage{lmodern}
\usepackage[margin=26mm,headheight=14pt]{geometry}
\usepackage{amsmath,amssymb,amsthm,mathtools}
\usepackage{microtype,booktabs,array,longtable,tabularx}
\usepackage{enumitem,xcolor,xurl,fancyhdr}
\usepackage{hyperref}
\usepackage[nameinlink,noabbrev]{cleveref}
\definecolor{Ink}{HTML}{253B4B}
\hypersetup{colorlinks=true,linkcolor=Ink,citecolor=Ink,urlcolor=Ink,
 pdftitle={Differential Surcomplex Numbers: Bounded Phase, Oscillation Obstructions, and Rank-One Equations},
 pdfauthor={Research article prepared for Vladimir Reshetnikov},
 pdfsubject={Internal differential algebra over the surreal and surcomplex fields}}
\pagestyle{fancy}\fancyhf{}
\fancyhead[L]{\small Differential surcomplex numbers}
\fancyhead[R]{\small Bounded phase and differential equations}
\fancyfoot[C]{\thepage}
\renewcommand{\headrulewidth}{0.3pt}
\setlength{\parskip}{0.22em}
\setlength{\emergencystretch}{2em}
\setlist{leftmargin=2em,itemsep=0.2em,topsep=0.35em}
\setcounter{tocdepth}{2}
\numberwithin{equation}{section}
\newtheorem{theorem}{Theorem}[section]
\newtheorem{proposition}[theorem]{Proposition}
\newtheorem{lemma}[theorem]{Lemma}
\newtheorem{corollary}[theorem]{Corollary}
\theoremstyle{definition}
\newtheorem{definition}[theorem]{Definition}
\newtheorem{example}[theorem]{Example}
\theoremstyle{remark}
\newtheorem{remark}[theorem]{Remark}
\crefname{theorem}{Theorem}{Theorems}
\crefname{proposition}{Proposition}{Propositions}
\crefname{lemma}{Lemma}{Lemmas}
\crefname{corollary}{Corollary}{Corollaries}
\crefname{definition}{Definition}{Definitions}
\crefname{example}{Example}{Examples}
\crefname{remark}{Remark}{Remarks}
\crefname{section}{Section}{Sections}
\newcommand{\No}{\mathbf{No}}
\newcommand{\F}{\mathbb F}
\newcommand{\E}{\mathbb E}
\newcommand{\R}{\mathbb R}
\newcommand{\C}{\mathbb C}
\newcommand{\Q}{\mathbb Q}
\newcommand{\Z}{\mathbb Z}
\newcommand{\N}{\mathbb N}
\newcommand{\OO}{\mathcal O}
\newcommand{\mm}{\mathfrak m}
\newcommand{\J}{\mathcal J}
\newcommand{\D}{\partial}
\newcommand{\I}{\mathcal I}
\newcommand{\Ex}{\operatorname{Exp}_{\mm}}
\newcommand{\Lo}{\operatorname{Log}_{\mm}}
\newcommand{\ob}{\operatorname{ob}}
\newcommand{\st}{\operatorname{st}}
\newcommand{\supp}{\operatorname{supp}}
\newcommand{\Ree}{\operatorname{Re}}
\newcommand{\Imm}{\operatorname{Im}}
\newcommand{\ld}{\operatorname{ld}}
\newcommand{\GL}{\operatorname{GL}}
\newcommand{\tr}{\operatorname{tr}}
\newcommand{\gm}{\mathbf G_{\mathrm m}}
\newcommand{\up}[1]{#1^{\uparrow}}
\newcommand{\down}[1]{#1^{\downarrow}}
\newcommand{\abs}[1]{\lvert#1\rvert}
\newcolumntype{Y}{>{\raggedright\arraybackslash}X}
\newcolumntype{P}[1]{>{\raggedright\arraybackslash}p{#1}}
\newcommand{\repobase}{https://github.com/VladimirReshetnikov/Surreal/blob/e260237db9b71da8b74a0c13c8e6355119091100/}

\begin{document}
\hypersetup{pageanchor=false}
\begin{titlepage}
\centering
\vspace*{1.0cm}
{\large\scshape A companion to the Surreal repository\par}
\vspace{0.8cm}
{\Huge\bfseries Differential\\[0.15em] Surcomplex Numbers\par}
\vspace{0.65cm}
{\Large Bounded Phase, Oscillation Obstructions,\\[0.2em]
and Rank-One Differential Equations\par}
\vspace{0.8cm}
{\large Research article prepared for Vladimir Reshetnikov\par}
\vspace{0.35cm}
{\large September 2026\par}
\vspace{0.9cm}
\begin{minipage}{0.94\textwidth}
\begin{abstract}
The repository develops a substantial theory of functions of surreal and
surcomplex variables, but treats the internal differential-field structure
mostly as a boundary condition on that theory. This article develops that
missing direction. Fix the simplest Berarducci--Mantova derivation on
$\F=\No$, normalized by $\D\omega=1$, and extend it to $\E=\No[i]$.
We prove that the logarithmic derivatives of nonzero surcomplex numbers are
exactly $\F+i\D\mm$, where $\mm$ is the real infinitesimal ideal. Consequently,
$\D y=ay$ has a nonzero surcomplex solution exactly when $\Imm a$ has a
\emph{finite real primitive}. The purely infinite part of that primitive is a
complete obstruction and a canonical invariant of rank-one equations under
gauge transformations.

This yields exact power and iterated-logarithm thresholds, a classification of
all homogeneous equations with ordinary complex constant coefficients, and an
existence criterion for fundamental matrices of upper-triangular systems.
For missing solutions we construct explicit Picard--Vessiot extensions and
prove that they introduce no new constants; integer relations among the
obstructions determine the differential Galois torus of a diagonal system.
We also explain why formal analytic solutions in an external variable do not
contradict these obstructions, and give a set-sized localization argument.
All principal deductions are proved. An accompanying standard-library Python
program passes 1,395 exact finite checks, which test examples and algebraic
identities rather than verify the infinitary theorems.
\end{abstract}
\end{minipage}
\vfill
\begin{minipage}{0.94\textwidth}\small
\textbf{Scope and status.} This is a mathematical exposition extending a
specific gap in a commit-pinned documentation snapshot. It does not claim
priority, a solution of a named published open problem, independent refereeing,
or proof-assistant verification. The Berarducci--Mantova existence and
integration theorems are imported from the published literature, not reproved.
The repository's existing phase decomposition and chosen-exponential
obstruction are explicitly acknowledged, not presented as missing results.
\end{minipage}
\end{titlepage}
\hypersetup{pageanchor=true}
\pagenumbering{roman}
\begingroup\small\setlength{\parskip}{0pt}\setlength{\baselineskip}{11pt}
\makeatletter
\renewcommand{\l@section}[2]{%
  \addpenalty{-\@highpenalty}\vskip0.45em
  \begingroup\bfseries\@dottedtocline{1}{0em}{2.1em}{#1}{#2}\endgroup}
\renewcommand{\l@subsection}{\@dottedtocline{2}{1.5em}{2.8em}}
\makeatother
\tableofcontents\endgroup
\clearpage\pagenumbering{arabic}

\section{The documentation gap and the main answer}\label{sec:gap}

\subsection{What the repository already does}
The audit concerns the repository \texttt{VladimirReshetnikov/Surreal} at
revision
\begin{center}\small\texttt{e260237db9b71da8b74a0c13c8e6355119091100}.\end{center}
Its documentation map lists fifteen research packages covering surreal
normal forms and special constructions, surcomplex analysis and geometry,
and foundations and computer algebra \cite{RepoMap}. The relevant gap is
\emph{not} the absence of differentiation. It is the absence of a developed
account of \emph{internal differential equations after complexifying the
surreal differential field} in the inspected central reports.

Three existing treatments delimit this gap particularly well. The computer
algebra report distinguishes differentiation in an external variable,
formal differentiation in a monomial parameter, and a derivation of the
surreal scalar field \cite[section with source label
\texttt{cas:sec-derivatives}]{RepoCAS}. The trigonometry report constructs
finite phases and explicitly says that its function derivatives are not
Berarducci--Mantova derivatives of scalars \cite{RepoTrig}. The analysis report
goes further: its ``Derivation obstruction'' theorem shows that its chosen
canonical global surcomplex exponential is incompatible with an internal
scalar derivation normalized by $\D\omega=1$ \cite[source label
\texttt{b:derivation-obstruction}]{RepoAnalysis}.

Thus there is already an obstruction theorem to build on. What is not
developed there is a criterion for \emph{every} first-order scalar equation,
a canonical classification under change of dependent variable, a systematic
treatment of the missing oscillatory solutions, or a differential Galois
interpretation. These are the central subjects of this article. The audit was
a targeted reading of the documentation map and the most relevant central
reports, not a claim that every sentence in every preserved source draft or
code file was examined.

The genetic-gap report is another useful boundary marker: its nonconstant
functions with pointwise derivative zero concern the fine order-field
derivative of a \emph{function}, not a scalar derivation \cite{RepoGaps}.
Neither its primitive nonuniqueness nor the analytic reports' many local
functions contradicts the constant-field theorem used below.

\subsection{The central result in accessible terms}
Fix the simplest Berarducci--Mantova derivation $\D$ and write
\[
 \F=\No,\qquad \E=\F[i],\qquad \D\omega=1.
\]
An equation such as $\D y=iy$ asks for a \emph{number} in the differential
field $\E$, not for a function of a new independent variable. The field
already contains every algebraic root of every finite-degree polynomial,
and every element has an additive primitive. Nonetheless, the equation
$\D y=iy$ has no nonzero solution in $\E$.

The reason is not a deficiency of algebraic closure. Every surcomplex unit
vector consists of a constant ordinary complex phase and an infinitesimal
phase. The scalar derivative kills the first and differentiates only the
second. A sustained nonzero ordinary angular frequency therefore cannot be
realized inside this field.

Here is the exact form of that principle. If $a=p+iq$ with $p,q\in\F$, then
\begin{equation}\label{eq:opening}
 \exists y\in\E^{\times}\ (\D y=ay)
 \quad\Longleftrightarrow\quad
 \exists B\in\F\ (\D B=q\ \text{and }B\text{ is finite}).
\end{equation}
The real coefficient $p$ never obstructs a solution. The issue is the
\emph{primitive} of the imaginary coefficient, not merely its magnitude.
For example,
\[
 \D y=\frac{i}{\omega}y \quad\text{has no nonzero solution},\qquad
 \D y=\frac{i}{\omega^2}y \quad\text{does}.
\]
Both coefficients are infinitesimal. Their primitives, $\log\omega$ and
$-1/\omega$, lie on opposite sides of the finite--infinite boundary.

\subsection{What is proved, and what is imported}
The proofs begin with an explicit differential-field contract in
\cref{sec:framework}. The existence of that structure on all surreals is the
deep published input. The subsequent phase calculation, obstruction map,
gauge classification, differential-equation results, and rational
no-new-constants arguments are supplied in full.

The development is not a new construction of the surreal field or its
exponential. Nor is it a claim that the elementary consequences here are
unknown in differential algebra. Classical Picard--Vessiot terminology and
linear differential-module methods are used in their usual sense
\cite{PS}; the purpose is to work out their precise content in this surreal
setting and connect them to the repository's analytic conventions.

\section{A precise differential-field framework}\label{sec:framework}

\subsection{The fixed derivation and its contract}
We use the \emph{simplest} Berarducci--Mantova derivation, not an arbitrary
map satisfying $\D\omega=1$. Its properties needed here are
\begin{align}
 \D(xy)&=x\D y+y\D x, & \ker(\D|_{\F})&=\R,\label{eq:contract1}\\
 \D\Bigl(\sum_{j\in S}x_j\Bigr)&=\sum_{j\in S}\D x_j,
 &\D(\exp x)&=(\exp x)\D x,\label{eq:contract2}\\
 \D\omega&=1, &\D\F&=\F,\qquad \D\mm\subseteq\mm.\label{eq:contract3}
\end{align}
The sum in \eqref{eq:contract2} is a set-indexed strongly summable Hahn
family. Here $\exp$ is the real surreal exponential and $\mm$ is the real infinitesimal ideal. These facts are
Theorems A and B, with the normalization stated between them, in
Berarducci--Mantova \cite{BM}; their construction and proof of surjectivity
are not reproduced here.

Throughout, ``finite'' means bounded in absolute value by an ordinary
positive integer. Put
\[
 \OO=\{x\in\F:\abs x<n\text{ for some }n\in\N\},\qquad
 \mm=\{x\in\F:\abs x<1/n\text{ for every }n\geq1\}.
\]
Every $x\in\OO$ is uniquely $r+\varepsilon$ with $r\in\R$ and
$\varepsilon\in\mm$. Consequently
\begin{equation}\label{eq:Dfinite}
 \D\OO=\D\mm\subseteq\mm.
\end{equation}
The inclusion is a statement about the chosen derivation. It is not a general
fact about all derivations of a non-Archimedean field.

\subsection{Three meanings of differentiation}
The following comparison should remain visible in any implementation.
\begin{center}\small
\begin{tabularx}{\textwidth}{@{}P{0.21\textwidth}YY@{}}
\toprule
Operator & What varies & Sample values\\
\midrule
External $D_z$ & A function's independent variable $z$; all surreal
coefficients are fixed & $D_z\omega=0$, $D_z z=1$\\[0.35em]
Formal $d/dt$ & A formal power or Hahn parameter & $(d/dt)t^3=3t^2$\\[0.35em]
Internal $\D$ & Elements of the chosen differential field &
$\D\omega=1$, $\D t=-t^2$ for $t=\omega^{-1}$\\
\bottomrule
\end{tabularx}
\end{center}
In particular, $\D(t^3)=-3t^4$. Replacing $d/dt$ by $\D$ in a displayed
ordinary differential equation without changing the coefficient functions
usually changes the equation.

The pointwise fine derivative of a class function is a fourth notion when
its topology and limit definition are part of the discussion. It is related
to $D_z$ on appropriate analytic functions, but it is not obtained by
applying $\D$ to the values while ignoring their dependence on the argument.
We will make a valid chain-rule comparison in \cref{sec:external}, rather
than identify these operators.

\subsection{Normal-form projections and a normalized primitive}
Split a Conway normal form at the monomial $1$:
\begin{equation}\label{eq:split}
 x=\up{x}+c_0(x)+\down{x},
 \qquad \up{x}\in\J,\quad c_0(x)\in\R,\quad\down{x}\in\mm.
\end{equation}
The class $\J$ consists of $0$ and all purely infinite normal forms, namely
those whose nonzero monomials are greater than $1$. These projections are
$\R$-linear; $c_0$ is not asserted to be multiplicative on all of $\F$.
We have
\[
 \F=\J\oplus\R\oplus\mm,\qquad \OO=\R\oplus\mm.
\]
In particular, $x$ is finite if and only if $\up{x}=0$.

\begin{definition}[Normalized primitive]\label[definition]{def:I}
For $b\in\F$, let $\I b$ be the unique real surreal satisfying
\[
 \D(\I b)=b,\qquad c_0(\I b)=0.
\]
For a complex input, put $\I_{\C}(p+iq)=\I p+i\I q$.
\end{definition}

\begin{proposition}\label[proposition]{prop:I}
The operator $\I$ exists, is $\R$-linear, and satisfies
\[
 \D\I b=b,\qquad \I\D x=x-c_0(x).
\]
Every primitive of $b$ differs from $\I b$ by a real constant.
\end{proposition}
\begin{proof}
Surjectivity gives a primitive $A$ of $b$. Subtracting $c_0(A)$ gives the
required normalization. Two such primitives differ by an element of
$\ker\D=\R$, and their constant normal-form coefficients force that
difference to be zero. Linearity and the two identities follow by the same
uniqueness argument.
\end{proof}
This is a semantic definition, not a universal finite algorithm for computing
primitives. It does not require a choice of one element from each primitive
family: the normalization singles out a unique element.

\subsection{Size conventions}
We may work in a class theory such as NBG with the displayed constructions
understood as class functions. Every individual surreal is a set-coded
object, every normal-form support is a set, and every polynomial, matrix,
and differential equation below has ordinary finite size. No proper-class
family of nonzero Hahn terms is summed.

There is an additional trap in differential classification. The notation
$\F/\OO$ would ordinarily mean a collection of cosets, but a coset of the
proper class $\OO$ is not an object that can serve as an element of another
class in NBG. We therefore use the \emph{representative class} $\J$ and the
map $x\mapsto\up{x}$ instead. References to a quotient mean the equivalence
relation represented this way, never a class whose members are proper
classes. Likewise, rank-one equations are classified by their scalar
coefficients and explicitly stated gauge relations, not by an illicit
collection of class-sized modules. Automorphism groups over a proper-class
base are encoded by the constant multipliers of finitely many generators;
we do not form a class whose members are proper-class maps. A set-sized
version is constructed in \cref{sec:localization}.

\section{Complexification, primitives, and infinitesimal calculus}\label{sec:complexification}

\subsection{The unique extension of the derivation}
\begin{proposition}\label[proposition]{prop:complexification}
There is a unique extension of $\D$ to $\E=\F[i]$, namely
\[
 \D(x+iy)=\D x+i\D y.
\]
It commutes with conjugation, its constant field is $\C$, and it is surjective.
\end{proposition}
\begin{proof}
Differentiating $i^2+1=0$ gives $2i\D i=0$, hence $\D i=0$. The displayed
formula is therefore necessary and directly satisfies the product rule.
It commutes with conjugation by inspection. If the derivative vanishes,
both real and imaginary parts are in $\R$, so the constant is in $\C$.
Finally $\I_{\C}$ supplies a primitive of every element.
\end{proof}

Algebraically, $\F$ is real closed and $\E$ is algebraically closed. It will
be important not to confuse these algebraic closure statements with the
existence of solutions to differential equations. Surjectivity of $\D$
solves $\D y=f$, but it does not by itself solve $\D y=ay$ nontrivially.

\begin{proposition}[Real exponential integration]\label[proposition]{prop:realexp}
Every $p\in\F$ is the logarithmic derivative of a positive element of $\F$.
More precisely, $s=\exp(\I p)$ satisfies $\D s=ps$.
\end{proposition}
\begin{proof}
Apply exponential compatibility to $\I p$.
\end{proof}
For real coefficients this is all that is needed to integrate first-order
homogeneous equations. The complex case requires more than appending the
symbol $i$ to this calculation: no globally differential-compatible complex
exponential has been provided.

\subsection{Strong summation near zero and one}
Write
\[
 \OO_{\C}=\OO+i\OO,\qquad \mm_{\C}=\mm+i\mm.
\]
For $h\in\mm_{\C}$ define
\begin{equation}\label{eq:small-exp-log}
 \Ex(h)=\sum_{n\geq0}\frac{h^n}{n!},\qquad
 \Lo(1+h)=\sum_{n\geq1}\frac{(-1)^{n+1}h^n}{n}.
\end{equation}
These are Hahn sums, not limits of finite partial sums in the full surreal
fine topology. The positive support condition is the relevant convergence
substitute: the support of an infinitesimal lies in the positive part of
the additive valuation group. The positive-support summability lemma ensures
that the monoid generated by such a well-ordered support is well ordered,
and that only finitely many words contribute to each fixed exponent.
Thus formal series identities at zero can be substituted in these
arguments. This is the same local mechanism used in the repository's
analytic and phase constructions \cite{RepoAnalysis,RepoTrig}.

\begin{lemma}[Infinitesimal exponential and logarithm]\label[lemma]{lem:small}
The maps
\[
 \Ex:(\mm_{\C},+)\longrightarrow(1+\mm_{\C},\cdot),\qquad
 \Lo:1+\mm_{\C}\longrightarrow\mm_{\C}
\]
are inverse group isomorphisms. They commute with conjugation and satisfy
\begin{align}
 \D\Ex(h)&=\Ex(h)\D h,\label{eq:Dsmall-exp}\\
 \D\Lo(1+h)&=\frac{\D h}{1+h}.\label{eq:Dsmall-log}
\end{align}
\end{lemma}
\begin{proof}
The formal identities for exponential and logarithm hold in characteristic
zero. Substitution at infinitesimals is valid by the support condition just
stated, including simultaneous substitution in finitely many infinitesimal
arguments. This proves inversion and the group laws. All scalar coefficients
are rational, so conjugation commutes with both maps.

For the derivative, the finite product rule gives
$\D(h^n)=nh^{n-1}\D h$. Strong additivity permits applying $\D$ termwise
in \eqref{eq:small-exp-log}. The resulting families are summable, and
factoring out the fixed element $\D h$ gives
\[
 \D\Ex(h)=\D h\sum_{n\geq1}\frac{h^{n-1}}{(n-1)!}.
\]
For the logarithm, the same calculation gives
$\D h\sum_{n\geq1}(-1)^{n+1}h^{n-1}=\D h/(1+h)$.
No exchange of analytic limits is being used.
\end{proof}

\section{Polar decomposition and the logarithmic-derivative image}\label{sec:polar}

\subsection{Every unit direction has an infinitesimal moving phase}
For $y=x+iz$ put $\bar y=x-iz$ and
$\abs y=(x^2+z^2)^{1/2}$. Standard parts of elements in $\OO_{\C}$ lie in
$\C$. If $u\bar u=1$, both coordinates of $u$ are finite and its standard
part belongs to the ordinary unit circle.

\begin{theorem}[Canonical polar data]\label[theorem]{thm:polar}
Every $y\in\E^{\times}$ has a unique representation
\begin{equation}\label{eq:polar}
 y=\exp(A)\,c\,\Ex(i\theta),
 \qquad A\in\F,\quad c\in\C,\quad c\bar c=1,\quad\theta\in\mm.
\end{equation}
The data are
\begin{equation}\label{eq:polar-data}
 A=\log\abs y,\qquad
 c=\st\!\left(\frac{y}{\abs y}\right),\qquad
 \theta=-i\Lo\!\left(\frac{y}{c\abs y}\right).
\end{equation}
The quantities $A$ and $\theta$ are additive under multiplication of $y$;
the quantity $c$ is multiplicative.
\end{theorem}
\begin{proof}
Put $r=\abs y>0$, $u=y/r$, and $c=\st(u)$. Then $c\bar c=1$ and
$w=u/c$ lies in $1+\mm_{\C}$ with $w\bar w=1$. By the group laws in
\cref{lem:small},
\[
 \overline{\Lo(w)}=\Lo(\bar w)=\Lo(w^{-1})=-\Lo(w).
\]
Hence $\Lo(w)$ is purely imaginary and $\theta=-i\Lo(w)$ belongs to the
real ideal $\mm$. The formula $w=\Ex(i\theta)$ proves existence.

Conversely, $\Ex(i\theta)$ has modulus one because its conjugate is its
inverse; its standard part is one. Therefore any representation of the
specified form must have $\exp(A)=\abs y$ and
$c=\st(y/\abs y)$. Injectivity of $\Ex$ on $\mm_{\C}$ determines
$\theta$. Multiplicativity of modulus and standard part, followed by the
logarithm group law, gives the last assertions.
\end{proof}

This is a precise version of the finite-angle decomposition already present
in the trigonometry report \cite{RepoTrig}. It avoids choosing an ordinary
argument branch: $c$ is a unit complex number rather than a selected real
angle. The important addition here is to apply the internal derivation to
these coordinates.

\subsection{The exact image of the logarithmic derivative}
\begin{definition}
For $y\in\E^{\times}$ write
$\ld(y)=\D y/y$. This is a homomorphism from the multiplicative group
$\E^{\times}$ to the additive group $\E$.
\end{definition}

\begin{theorem}[Logarithmic-derivative image]\label[theorem]{thm:image}
For the polar data in \cref{thm:polar},
\begin{equation}\label{eq:logder-polar}
 \ld(y)=\D A+i\D\theta.
\end{equation}
Consequently,
\begin{equation}\label{eq:image}
 \ld(\E^{\times})=\F+i\D\mm.
\end{equation}
The kernel of $\ld$ is $\C^{\times}$, and
$\D:\mm\to\D\mm$ is injective.
\end{theorem}
\begin{proof}
Differentiate \eqref{eq:polar}. The ordinary complex factor $c$ is constant,
the real exponential contributes $\D A$, and \cref{lem:small} gives
$i\D\theta$ for the final factor. This proves one inclusion in
\eqref{eq:image}. For the other, take any $p\in\F$ and any
$q=\D\theta$ with $\theta\in\mm$. Then
\[
 y=\exp(\I p)\Ex(i\theta)
\]
has logarithmic derivative $p+iq$. The kernel assertion follows from the
constant-field statement. Finally, $\mm\cap\R=\{0\}$, so $\D$ has
zero kernel on $\mm$.
\end{proof}

The product $\exp(\I p)\Ex(i\theta)$ uses two different exponential
constructions with explicitly different domains. It is not an abbreviation
for an unspecified global surcomplex exponential.

\begin{remark}[Small is not small enough]\label[remark]{rem:smallnotenough}
The inclusion $\D\mm\subseteq\mm$ is strict. Indeed,
$\D\log\omega=1/\omega$, and every primitive of $1/\omega$ differs from
$\log\omega$ by a real constant. Since $\log\omega$ is infinite,
$1/\omega\notin\D\mm$. Thus testing only whether the imaginary coefficient
is infinitesimal loses the decisive information.
\end{remark}

\section{The bounded-phase theorem and canonical gauge forms}\label{sec:obstruction}

\subsection{A complete solvability criterion}
\begin{theorem}[Bounded-phase criterion]\label[theorem]{thm:bounded}
Let $a=p+iq\in\E$. The following are equivalent:
\begin{enumerate}[label=(\roman*)]
\item $\D y=ay$ has a nonzero solution in $\E$;
\item $q\in\D\mm$;
\item $q$ has a finite real primitive;
\item $\I q$ is finite;
\item $\up{(\I q)}=0$.
\end{enumerate}
When these conditions hold, there is a unique $\theta\in\mm$ with
$\D\theta=q$. With $A=\I p$, all nonzero solutions are
\begin{equation}\label{eq:all-rank-one}
 y=c\exp(A)\Ex(i\theta),\qquad c\in\C^{\times}.
\end{equation}
\end{theorem}
\begin{proof}
The equivalence of (i) and (ii) is \cref{thm:image}. If $q=\D\theta$ with
$\theta\in\mm$, it has a finite primitive. Conversely, subtract the
standard part from any finite primitive to obtain an infinitesimal one.
All primitives differ by a real constant, proving (iii)$\Leftrightarrow$(iv).
The normal-form splitting proves (iv)$\Leftrightarrow$(v). Injectivity of
$\D$ on $\mm$ gives uniqueness of $\theta$, and the displayed expression
is a solution. If $y_1,y_2$ are nonzero solutions, the quotient rule gives
$\D(y_1/y_2)=0$, so their ratio is a nonzero complex constant.
\end{proof}

\subsection{A canonical obstruction map}
\begin{definition}\label[definition]{def:ob}
For $a\in\E$, define
\begin{equation}\label{eq:ob}
 \ob(a)=\up{(\I\Imm a)}\in\J.
\end{equation}
\end{definition}
The normalization of $\I$ is convenient but not essential here: adding a
real constant to a primitive does not change its purely infinite part.

\begin{theorem}[Exact obstruction sequence]\label[theorem]{thm:exact}
The map $\ob$ is $\R$-linear and surjective. Its kernel is
$\ld(\E^{\times})$. Thus the sequence
\begin{equation}\label{eq:exact}
 1\longrightarrow\C^{\times}\longrightarrow\E^{\times}
 \xrightarrow{\ \ld\ }(\E,+)
 \xrightarrow{\ \ob\ }(\J,+)\longrightarrow0
\end{equation}
is exact, with multiplicative group laws before $\ld$ and additive group
laws after it. Equivalently, $\J$ gives concrete representatives for the
obstruction quotient.
\end{theorem}
\begin{proof}
Linearity follows from the linearity of the imaginary-part, primitive,
and normal-form projection maps. The kernel statement is
\cref{thm:bounded}. If $B\in\J$, then $c_0(B)=0$ and
$\I\D B=B$, so $\ob(i\D B)=B$. This proves surjectivity. Exactness at
$\E^{\times}$ is the kernel assertion of \cref{thm:image}.
\end{proof}

Although $\E$ is a complex field, $\ob$ is not a complex-linear map. In
particular, all real coefficients have zero obstruction, whereas multiplying
one of them by $i$ may produce a nonzero obstruction. Its linearity over
$\R$ will be useful; its integer linearity controls differential Galois
relations later.

\subsection{Gauge transformations and uniqueness of the normal form}
A change of dependent variable $y=gz$, with $g\in\E^{\times}$, sends
\begin{equation}\label{eq:gauge}
 \D y=ay\qquad\text{to}\qquad
 \D z=\bigl(a-\ld(g)\bigr)z.
\end{equation}
This is a \emph{gauge transformation}. It does not change whether a
nonzero solution exists in the base field.

\begin{theorem}[Purely infinite phase normal form]\label[theorem]{thm:gauge}
Every equation $\D y=ay$ is gauge-equivalent to exactly one equation of
the form
\begin{equation}\label{eq:normal}
 \D z=i\D B\,z,\qquad B\in\J.
\end{equation}
The canonical parameter is $B=\ob(a)$. Two coefficients $a,a'\in\E$ are
gauge-equivalent if and only if $\ob(a)=\ob(a')$.
\end{theorem}
\begin{proof}
Write $a=p+iq$, $A=\I p$, and $\I q=B+\varepsilon$, where
$B\in\J$ and $\varepsilon\in\mm$. There is no constant term because
$\I q$ is normalized. Take
\[
 g=\exp(A)\Ex(i\varepsilon).
\]
Then $\ld(g)=p+i\D\varepsilon$, so
$a-\ld(g)=i\D B$. This proves existence.

Gauge-equivalent coefficients differ by an element of $\ker\ob$, and
therefore have the same obstruction. Conversely, if their obstructions
agree, their difference is a logarithmic derivative by \cref{thm:exact},
which supplies a gauge transformation. Finally,
$\ob(i\D B)=B$, proving uniqueness of the representative.
\end{proof}

The parameter $B$ is not an angular value assigned to an infinite input.
It records precisely the part of a primitive that \emph{cannot} be absorbed
by an existing surcomplex phase. The actual solution factor absorbed by the
gauge transformation has only the infinitesimal phase $\varepsilon$.

\begin{corollary}[Rank-one tensor calculus]\label[corollary]{cor:rankone}
The gauge classes of first-order homogeneous equations, with tensor product
represented by addition of their coefficients, are represented by the
additive class $\J$. Duality corresponds to $B\mapsto-B$. Every nonzero
class has infinite order.
\end{corollary}
\begin{proof}
The coefficient of the tensor product is $a+a'$; that of the dual is
$-a$. Apply the linearity of $\ob$. Since $\J$ is a real vector space,
$nB=0$ for a nonzero integer $n$ implies $B=0$.
\end{proof}
This rank-one differential classification is different from the holomorphic
line-bundle or divisor classifications in the repository's global analytic
reports. The equivalence relation here is differential gauge equivalence.

\section{Oscillation obstructions and exact thresholds}\label{sec:examples}

\subsection{The basic missing oscillator}
\begin{corollary}\label[corollary]{cor:noosc}
The equation $\D y=iy$ has no nonzero solution in $\E$.
More generally, $\D y=\lambda y$ with $\lambda\in\C$ has a nonzero
solution in $\E$ if and only if $\lambda\in\R$.
\end{corollary}
\begin{proof}
If $\lambda=r+is$, then $\I(\Imm\lambda)=s\omega$.
This is finite exactly when $s=0$. Apply \cref{thm:bounded}.
\end{proof}
There is a short independent proof of the first assertion that makes the
role of smallness especially transparent. Suppose $\D y=iy$ with $y\ne0$.
Then
\[
 \D(y\bar y)=iy\bar y-iy\bar y=0.
\]
Thus $r=\abs y$ is a positive real constant. The unit $u=y/r$ satisfies
$\D u=iu$. Its two coordinates are finite, so their derivatives are
infinitesimal by \eqref{eq:Dfinite}. The equation forces both coordinates
of $u$ to be infinitesimal, contradicting $u\bar u=1$.

This argument uses an internal differential field, not an assertion that
ordinary sine or cosine fails to exist as a function. Nor does it invalidate
finite surreal angles. It says that an element with this sustained internal
angular velocity cannot lie in $\No[i]$.

\subsection{A power-scale Hahn field}
Let $t=\omega^{-1}$. Restrict for the moment to the real-exponent Hahn field
\[
 H_{\R}=\R((t^{\R})),\qquad H_{\C}=\C((t^{\R})),
\]
embedded in the surreal and surcomplex fields by Conway monomials.
For an ordinary real exponent $\alpha$,
\begin{equation}\label{eq:power-D}
 \D(t^{\alpha})=-\alpha t^{\alpha+1}.
\end{equation}
Here real powers agree with the real exponential definition. This formula
must not be extended to arbitrary surreal exponents: Conway's monomial
map and general surreal exponentiation are different constructions. The
normal-form and exponential conventions are those of \cite{BM}.
On this specified Hahn field,
\[
 \D=-t^2\frac{d}{dt}.
\]

\begin{theorem}[Sharp power-scale criterion]\label[theorem]{thm:power}
Let
$ b=\sum_{\alpha\in S}b_\alpha t^\alpha\in H_{\R}$, where $S$ is
well ordered, and put absent coefficients equal to zero. A primitive in
$\F$ is
\begin{equation}\label{eq:powerprimitive}
 B=\sum_{\substack{\alpha\in S\\\alpha\ne1}}
       \frac{b_\alpha}{1-\alpha}\,t^{\alpha-1}
       +b_1\log\omega.
\end{equation}
For every $p\in\F$, the equation $\D y=(p+ib)y$ has a nonzero solution
in $\E$ if and only if $b=0$ or
\begin{equation}\label{eq:powerthreshold}
 \min\supp(b)>1.
\end{equation}
In the soluble case, the sum in \eqref{eq:powerprimitive} is the unique
infinitesimal primitive of $b$.
\end{theorem}
\begin{proof}
Shifting the well-ordered support by $-1$ preserves well ordering, and
omitting one exponent and adding the one logarithmic term preserves
summability. Formula \eqref{eq:power-D} and strong additivity give
$\D B=b$. If every exponent is greater than $1$, all powers in the
primitive have positive exponents, there is no logarithmic term, and $B$
is infinitesimal.

If the least exponent $\alpha_0$ is less than $1$, its primitive term has
negative exponent and nonzero coefficient. It dominates the remaining
power terms and also $\log\omega$: for every real $\delta>0$,
$\log\omega\prec\omega^{\delta}$. Thus $B$ is infinite. If the least
exponent is $1$, the logarithmic term has nonzero coefficient while all
remaining power terms are infinitesimal, so $B$ is again infinite. No real
constant added to $B$ changes either conclusion. Apply
\cref{thm:bounded}.
\end{proof}
The elementary logarithmic comparison used here is a standard consequence
of the surreal real exponential; equivalently, it follows from the
normal-form description of $\log\omega$ \cite{BM}. The result is a
\emph{strict} threshold at $1$, not just a positivity-of-valuation condition.

\begin{example}[One scale, four outcomes]
The following table uses $\Ex$ only on infinitesimal arguments.
\begin{center}\small
\begin{tabularx}{\textwidth}{@{}P{0.25\textwidth}P{0.24\textwidth}Y@{}}
\toprule
Coefficient $a$ & Primitive of $\Imm a$ & Nonzero solution or obstruction\\
\midrule
$i/\omega$ & $\log\omega$ & No nonzero solution\\[0.3em]
$i/\omega^2$ & $-1/\omega$ & $\Ex(-i/\omega)$\\[0.3em]
$1+i/\omega^2$ & $-1/\omega$ & $\exp(\omega)\Ex(-i/\omega)$\\[0.3em]
$1+i/\omega$ & $\log\omega$ & No nonzero solution\\
\bottomrule
\end{tabularx}
\end{center}
A real growth factor can be arbitrarily large without repairing an infinite
phase obstruction.
\end{example}

\subsection{Why the threshold cannot be globalized by valuation alone}
The preceding theorem is deliberately restricted to real-exponent Hahn
coefficients. Put $L=\log\omega$. Then
\[
 \D(\log L)=\frac{1}{\omega L},\qquad
 \D(-L^{-1})=\frac{1}{\omega L^2}.
\]
The first primitive is infinite and the second infinitesimal. Therefore
\begin{equation}\label{eq:log-separator}
 \D y=\frac{i}{\omega\log\omega}y
 \ \text{has no nonzero solution},\qquad
 \D y=\frac{i}{\omega(\log\omega)^2}y
 \ \text{has a nonzero solution}.
\end{equation}
In the full Hahn valuation normalized by $v(t)=1$, one has $v(L)<0$, so
\[
 v(t/L)=1-v(L)>1.
\]
Thus even $v(\Imm a)>1$ is not a sufficient criterion over all of $\F$.
The missing information lies in the slower logarithmic scale. A CAS must
not apply \cref{thm:power} outside its declared coefficient domain.

\begin{theorem}[Iterated-logarithm threshold]\label[theorem]{thm:logs}
For ordinary integers $n\geq0$, define
$L_0=\omega$ and $L_{n+1}=\log L_n$. For a real number $p$, let
\[
 b_{n,p}=\frac{1}{L_0L_1\cdots L_{n-1}L_n^{p}},
\]
with the preceding product empty when $n=0$. Then
$\D y=i b_{n,p}y$ has a nonzero solution in $\E$ exactly when $p>1$.
\end{theorem}
\begin{proof}
Every $L_n$ is positive infinite. The exponential chain rule gives
\[
 \D L_n=\frac{1}{L_0\cdots L_{n-1}}.
\]
A primitive of $b_{n,p}$ is
\[
 B_{n,p}=
 \begin{cases}
  L_n^{1-p}/(1-p),&p\ne1,\\
  L_{n+1},&p=1.
 \end{cases}
\]
This is infinitesimal for $p>1$ and infinite for $p\leq1$.
The bounded-phase criterion proves the claim.
\end{proof}
There are therefore successively finer versions of the same boundary.
A primitive-based representation handles them uniformly; a single power
cutoff does not.

\subsection{A rational unit-circle example}
Consider
\[
 u=\frac{1+it}{1-it},\qquad t=\omega^{-1}.
\]
It has unit modulus and standard part $1$. Direct algebra gives
\begin{equation}\label{eq:cayley}
 \ld(u)=\frac{-2it^2}{1+t^2}.
\end{equation}
Its infinitesimal phase is $2\arctan(t)$, where the inverse tangent is the
infinitesimal Taylor evaluation. Indeed,
\[
 \D\bigl(2\arctan(t)\bigr)=\frac{-2t^2}{1+t^2}.
\]
This example shows that bounded-phase solutions need not be stored as
infinite sums: an exact rational expression can already carry the required
phase. Its full phase expansion is optional.

\section{All ordinary constant-coefficient homogeneous equations}\label{sec:constant}

Complex algebraic closure does not give the usual exponential-polynomial
solution space for internal differential equations. Precisely the real
exponential modes survive.

\begin{lemma}\label[lemma]{lem:Dpower}
For every ordinary integer $m\geq1$,
\[
 \ker(\D^m)=\{P(\omega):P\in\C[X],\ \deg P<m\}.
\]
\end{lemma}
\begin{proof}
For $m=1$ this is the constant-field theorem. Suppose the result holds for
$m-1$ and $\D^m y=0$. Then $\D y$ is a complex polynomial in $\omega$
of degree less than $m-1$. Integrate that polynomial term by term using
$\D\omega=1$. The difference between $y$ and the resulting polynomial
has derivative zero and is a complex constant. The converse follows by
finite differentiation.
\end{proof}

\begin{theorem}[The surviving constant-coefficient modes]\label[theorem]{thm:constant}
Let $P(T)\in\C[T]$ be nonzero, and let $m_\lambda$ be the multiplicity
of a root $\lambda$ of $P$. Then
\begin{equation}\label{eq:constant-kernel}
 \ker P(\D)=
 \bigoplus_{\substack{\lambda\in\R\\P(\lambda)=0}}
 \exp(\lambda\omega)
 \{Q(\omega):Q\in\C[X],\ \deg Q<m_\lambda\}.
\end{equation}
Its dimension over $\C$ is the sum of the multiplicities of the real roots
of $P$. In particular, an operator of positive order may have zero kernel.
\end{theorem}
\begin{proof}
For real $\lambda$, exponential compatibility gives
\[
 (\D-\lambda)\bigl(\exp(\lambda\omega)z\bigr)
       =\exp(\lambda\omega)\D z.
\]
Iterating and using \cref{lem:Dpower} identifies the kernel of
$(\D-\lambda)^{m_\lambda}$.

If $\lambda$ is nonreal, \cref{cor:noosc} says that $\D-\lambda$ is
injective. Every positive power is then injective as well.

It remains to combine the factors. Distinct factors
$(T-\lambda)^{m_\lambda}$ are pairwise coprime in $\C[T]$. The usual
B\'ezout identities for these factors, evaluated at the $\C$-linear
operator $\D$, decompose the kernel of their product into the direct sum
of their kernels. This is a finite polynomial identity and does not require
any topology or completeness. Only the real-root summands are nonzero,
giving \eqref{eq:constant-kernel}. Their displayed polynomial bases give
the dimension.
\end{proof}

\begin{example}
The following are identities of solution spaces inside $\E$:
\begin{align*}
 \D^2y+y=0&\quad\Longrightarrow\quad y=0,\\
 \D^2y-y=0&\quad\Longleftrightarrow\quad
 y=c_1\exp\omega+c_2\exp(-\omega),\\
 (\D-1)^2(\D^2+1)y=0&\quad\Longleftrightarrow\quad
 y=(c_1+c_2\omega)\exp\omega,
 \qquad c_1,c_2\in\C.
\end{align*}
The missing modes in the first and third equations are precisely the
oscillatory ones.
\end{example}

The coefficient hypothesis matters: the coefficients of $P$ are ordinary
complex constants. For a coefficient $a\in\E\setminus\C$, the operators
$\D$ and multiplication by $a$ do not commute. Polynomial factorization in
a commuting symbol cannot then be transported blindly to differential
operators.

\begin{corollary}[A Riccati equation]\label[corollary]{cor:riccati}
The only solutions in $\E$ of
\[
 \D u=1+u^2
\]
are $u=i$ and $u=-i$. In particular, there is no solution in $\F$.
\end{corollary}
\begin{proof}
The two stated constants are solutions. For any other solution,
$v=(u-i)/(u+i)$ is defined and nonzero, and
\[
 \ld(v)=\frac{\D u}{u-i}-\frac{\D u}{u+i}
       =\frac{2i\D u}{u^2+1}=2i.
\]
This contradicts \cref{cor:noosc}.
\end{proof}
The example is not a general nonlinear solvability theorem. It is a direct
way to transfer the first-order phase obstruction to a familiar nonlinear
equation.

\section{Inhomogeneous equations and triangular systems}\label{sec:systems}

\subsection{Variation of constants, with the obstruction visible}
\begin{proposition}\label[proposition]{prop:variation}
Suppose $a,f\in\E$ and $\ob(a)=0$. Choose a nonzero $s\in\E$ with
$\D s=as$. Then all solutions of
\[
 \D y-ay=f
\]
are
\begin{equation}\label{eq:variation}
 y=s\bigl(\I_{\C}(f/s)+c\bigr),\qquad c\in\C.
\end{equation}
If $\ob(a)\ne0$, the inhomogeneous equation has at most one solution in
$\E$.
\end{proposition}
\begin{proof}
Write $y=sz$. The equation becomes $\D z=f/s$, whose solutions are the
normalized primitive plus a complex constant. If the homogeneous equation
has no nonzero solution, the difference of two inhomogeneous solutions
must be zero.
\end{proof}

``At most one'' must not be replaced by ``none''. For example,
\[
 (\D-i)i=1,
 \qquad (\D-i)(i\omega+1)=\omega.
\]
Both equations have unique solutions in $\E$, despite the absence of a
nonzero solution to the corresponding homogeneous equation.

\begin{proposition}[Polynomial forcing]\label[proposition]{prop:polyforcing}
For $f\in\C[\omega]$ and $\lambda\in\C^{\times}$, a solution of
$(\D-\lambda)y=f$ is
\begin{equation}\label{eq:polyforcing}
 y=-\sum_{k=0}^{\deg f}\frac{\D^k f}{\lambda^{k+1}}.
\end{equation}
For $f=0$, take the displayed particular solution to be zero. If
$\lambda$ is nonreal this is the unique solution in $\E$; if $\lambda$
is real, all solutions are obtained by adding
$c\exp(\lambda\omega)$, $c\in\C$.
\end{proposition}
\begin{proof}
Applying $\D-\lambda$ to the finite sum makes all intermediate terms
cancel. The last derivative $\D^{\deg f+1}f$ vanishes, leaving $f$.
The homogeneous classification gives the rest.
\end{proof}

These distinctions also prevent a terminology error. A statement that an
operator has trivial kernel is not a statement that it fails to be
surjective. In particular, the preceding results do not claim to settle
surjectivity of all linear differential operators on $\E$.

\subsection{Fundamental matrices for upper-triangular systems}
Let $A=(a_{ij})$ be an upper-triangular $n\times n$ matrix over $\E$,
with $n$ an ordinary positive integer. A fundamental matrix is
$Y\in\GL_n(\E)$ satisfying $\D Y=AY$, with entrywise differentiation.

\begin{theorem}[Triangular fundamental-matrix criterion]\label[theorem]{thm:triangular}
There is a fundamental matrix in $\GL_n(\E)$ if and only if
\begin{equation}\label{eq:diagcriterion}
 \ob(a_{jj})=0\qquad(1\leq j\leq n).
\end{equation}
When the criterion holds, an upper-triangular fundamental matrix can be
constructed by finitely many applications of first-order homogeneous
integration and additive primitives.
\end{theorem}
\begin{proof}
For sufficiency, choose $s_j\ne0$ with
$\D s_j=a_{jj}s_j$, set $Y_{jj}=s_j$, and set $Y_{ij}=0$ for $i>j$.
For each column $j$, construct the entries above the diagonal in descending
order of $i$ by
\begin{equation}\label{eq:triangular-recursion}
 Y_{ij}=s_i\I_{\C}\!\left(
 s_i^{-1}\sum_{k=i+1}^{j}a_{ik}Y_{kj}\right),\qquad i<j.
\end{equation}
The entries on the right have already been constructed. The product rule
verifies
$\D Y_{ij}=a_{ii}Y_{ij}+\sum_{k=i+1}^{j}a_{ik}Y_{kj}$.
Thus $\D Y=AY$, and $\det Y=\prod_j s_j\ne0$.

For necessity, use induction on $n$. The last row of any fundamental
matrix satisfies $\D Y_{nj}=a_{nn}Y_{nj}$. At least one entry is nonzero,
so $\ob(a_{nn})=0$. All the last-row entries are complex constant
multiples of one nonzero solution. A change of columns by an invertible
constant matrix therefore makes the last row $(0,\ldots,0,s_n)$.
The resulting matrix remains fundamental and has block form
\[
 \begin{pmatrix}Z&v\\0&s_n\end{pmatrix},\qquad \det Z\ne0.
\]
The upper-left block satisfies the system for the leading
$(n-1)\times(n-1)$ block of $A$. Induction proves the remaining diagonal
conditions.
\end{proof}

In particular, every strictly upper-triangular system has a unipotent
fundamental matrix over $\E$. This is a finite iterated-integration fact:
there is no need to invoke an infinite exponential series for an arbitrary
matrix.

\subsection{The determinant test is necessary but insufficient}
For any finite matrix system, not necessarily triangular, Jacobi's formula
gives
\[
 \D\det Y=(\tr A)\det Y
\]
whenever $Y$ is fundamental. Hence
\begin{equation}\label{eq:tracenecessary}
 \ob(\tr A)=0
\end{equation}
is necessary. It is not sufficient. For
\[
 A=\begin{pmatrix}i&0\\0&-i\end{pmatrix},
\]
the trace vanishes, but neither diagonal equation has a nonzero solution.
Indeed, the vector system has no nonzero solution at all. Opposite phase
obstructions can cancel in a determinant without disappearing from the
individual modes. A determinant-only test therefore discards essential
differential information.

\section{Adjoining the missing solutions}\label{sec:PV}

\subsection{A rank-one Picard--Vessiot extension}
The obstruction says which solutions are absent from $\E$. It also gives
a direct construction of the extension needed to add one. We use
Picard--Vessiot terminology as in \cite{PS}: a solution-generated
differential extension with no new constants. All arguments are finite
rational-function arguments. For the proper-class base they can be read
as explicit class constructions; \cref{sec:localization} supplies an
ordinary set-sized setting whenever that is desired.

\begin{theorem}[The missing-solution extension]\label[theorem]{thm:PVone}
Suppose $a\in\E$ and $\ob(a)\ne0$. Introduce a transcendental symbol $U$
and extend the derivation to $\E(U)$ by
\begin{equation}\label{eq:DU}
 \D U=aU.
\end{equation}
Then the constant field of $\E(U)$ is exactly $\C$. The extension is a
Picard--Vessiot extension for $\D y=ay$, and its differential automorphisms
over $\E$ are exactly
\[
 U\longmapsto cU,\qquad c\in\C^{\times}.
\]
Thus its differential Galois group is $\gm(\C)$. If $\ob(a)=0$, the
corresponding Picard--Vessiot extension is trivial.
\end{theorem}
\begin{proof}
The product and quotient rules uniquely extend the derivation from
$\E[U]$ to $\E(U)$. The main point is to rule out new constants.
Let $R=P(U)/Q(U)\ne0$ have derivative zero, with nonzero coprime
$P,Q\in\E[U]$. Then
\begin{equation}\label{eq:rationalconstant}
 Q\D P=P\D Q.
\end{equation}
Coprimality implies $P\mid\D P$ and $Q\mid\D Q$.
Since $\D U=aU$, differentiation does not increase polynomial degree in
$U$. Therefore
\[
 \D P=hP,\qquad \D Q=hQ
\]
for the same $h\in\E$; equality follows from
\eqref{eq:rationalconstant}.

Suppose $p_jU^j$ and $p_kU^k$ are two nonzero terms of $P$ with $j\ne k$.
Comparing coefficients gives
\[
 \frac{\D p_j}{p_j}+ja=h
   =\frac{\D p_k}{p_k}+ka,
\]
so
\begin{equation}\label{eq:coefficient-relation}
 \ld(p_j/p_k)=(k-j)a.
\end{equation}
Applying $\ob$ yields $(k-j)\ob(a)=0$, impossible because $\J$ is
 torsion-free. Hence $P$ has just one term. The same argument applies to
$Q$, so $R=cU^m$ for some $c\in\E^{\times}$ and $m\in\Z$. The
constant equation is $\ld(c)+ma=0$. Again $\ob$ forces $m=0$, and then
$\D c=0$, giving $c\in\C^{\times}$.

This proves that the extension has no new constants. It is generated by
a nonzero fundamental solution and its inverse, as required. Any
differential automorphism sends $U$ to another nonzero solution; their
ratio is a constant, so it has the asserted form. Conversely, every such
substitution is a differential automorphism. When the obstruction is zero,
a nonzero solution already lies in $\E$ by \cref{thm:bounded}.
\end{proof}

The transcendence here is necessary: algebraic closure of $\E$ leaves no
nontrivial algebraic extension in which a missing solution could appear.
Moreover, taking an algebraic root cannot cancel a nonzero phase
obstruction, since every nonzero integer multiple of that obstruction is
still nonzero. The new symbol is a genuinely new differential element,
not an alternative notation for a pre-existing surcomplex number.

\subsection{A formal oscillator and its real form}
For $a=i$, extend conjugation to $\E(U)$ by $\bar U=U^{-1}$. This is
compatible with the derivation because
\[
 \D(U^{-1})=-iU^{-1}=\overline{\D U}.
\]
Define
\begin{equation}\label{eq:formaloscillator}
 c_{\mathrm{osc}}=\frac{U+U^{-1}}2,
 \qquad
 s_{\mathrm{osc}}=\frac{U-U^{-1}}{2i}.
\end{equation}
These elements are fixed by conjugation and satisfy
\[
 \D c_{\mathrm{osc}}=-s_{\mathrm{osc}},\qquad
 \D s_{\mathrm{osc}}=c_{\mathrm{osc}},\qquad
 c_{\mathrm{osc}}^2+s_{\mathrm{osc}}^2=1.
\]
Thus a formal oscillator is perfectly consistent in a larger differential
field. It is inconsistent only with insisting that its coordinates already
belong to the specified surreal differential field.

A rational real presentation makes the distinction concrete. Introduce a
transcendental $T$ over $\F$ with
\[
 \D T=\frac{1+T^2}{2},\qquad
 c_{\mathrm{osc}}=\frac{1-T^2}{1+T^2},\qquad
 s_{\mathrm{osc}}=\frac{2T}{1+T^2}.
\]
The same three identities follow by the quotient rule. The two
presentations are related by $U=(1+iT)/(1-iT)$, so the real rational
field has constant field $\R$, by \cref{thm:PVone}. The rational
function field $\F(T)$ admits field orders extending the order on $\F$,
but no such realization can preserve the derivative-smallness property
$\D\OO\subseteq\mm$ with constants $\R$: both oscillator coordinates
would be bounded by one, their derivatives would be infinitesimal, and the
unit-circle identity would give the same contradiction as before.

Consequently these are not definitions of $\sin\omega$ and $\cos\omega$
as BM-differential elements of $\No$. They are coordinates of a new formal
phase in a larger field whose differential-order behavior has changed.

\section{Several phases and their differential Galois torus}\label{sec:torus}

\subsection{Integer relations are the right relations}
Let $a_1,\ldots,a_m\in\E$, and consider the diagonal system
\[
 \D y_j=a_jy_j\qquad(1\leq j\leq m).
\]
Write $B_j=\ob(a_j)\in\J$. Define its phase-relation lattice
\begin{equation}\label{eq:lattice}
 L=\left\{n=(n_1,\ldots,n_m)\in\Z^m:
               \sum_{j=1}^{m}n_jB_j=0\right\}.
\end{equation}
If $n\in L$, then $\sum n_ja_j$ is a logarithmic derivative in $\E$.
Correspondingly, a product of powers of solutions can be normalized to an
element of the base field. These are integer relations because rational
functions in solutions involve integral monomial exponents. Real linear
dependence of phases is not enough.

\begin{lemma}\label[lemma]{lem:saturated}
The lattice $L$ is saturated in $\Z^m$: if $kn\in L$ for a nonzero
integer $k$, then $n\in L$. Hence $\Z^m/L$ is a free abelian group of
rank $d=m-\operatorname{rank}L$.
\end{lemma}
\begin{proof}
If $k\sum n_jB_j=0$, torsion-freeness of $\J$ gives
$\sum n_jB_j=0$. Thus the quotient is torsion-free and finitely generated,
so it is free. Equivalently, the elementary integer normal-form theorem
extends a basis of this saturated subgroup to a basis of $\Z^m$.
\end{proof}

\subsection{Construction and no-new-constants proof}
\begin{theorem}[The diagonal phase torus]\label[theorem]{thm:torus}
The diagonal system has a Picard--Vessiot extension obtained by adjoining
$d$ algebraically independent solution symbols after base-field gauges and
an invertible integer monomial change of coordinates. Its differential
Galois group is isomorphic to $\gm(\C)^d$. In the original coordinates,
its action on the solutions is described by
\begin{equation}\label{eq:torus}
 G=\left\{(c_1,\ldots,c_m)\in(\C^{\times})^m:
       \prod_{j=1}^{m}c_j^{n_j}=1\quad\text{for every }n\in L\right\}.
\end{equation}
\end{theorem}
\begin{proof}
Choose an integer basis of $\Z^m$ whose first
$\ell=\operatorname{rank}L$ vectors form a basis of $L$. The associated
unimodular matrix changes the diagonal coefficients to integer linear
combinations $a'_1,\ldots,a'_m$. The first $\ell$ combinations have zero
obstruction and therefore have nonzero solutions $g_j\in\E$.
The remaining obstructions $B'_{\ell+1},\ldots,B'_m$ are linearly
independent over $\Z$: any relation among them would be an element of
$L$ with nonzero complementary coordinates.

For these $d$ remaining coefficients introduce algebraically independent
symbols $U_1,\ldots,U_d$ with
$\D U_k=a'_{\ell+k}U_k$. We prove that the rational extension
$\E(U_1,\ldots,U_d)$ has no new constants. If a nonzero rational constant
is written as $P/Q$ with coprime multivariate polynomials, the same
coprimality argument as in \cref{thm:PVone} yields
$\D P=hP$ and $\D Q=hQ$ for $h\in\E$. Indeed, total degree does not
increase under this derivation, so each polynomial quotient is a scalar.

If $p_\alpha U^\alpha$ and $p_\beta U^\beta$ are distinct nonzero terms
of $P$, coefficient comparison gives
\[
 \ld(p_\alpha/p_\beta)
    =\sum_{k=1}^{d}(\beta_k-\alpha_k)a'_{\ell+k}.
\]
Applying $\ob$ contradicts integer independence unless $\alpha=\beta$.
Thus both $P$ and $Q$ are monomials. Applying the same argument to their
quotient shows that a rational constant has zero exponent vector and a
coefficient in $\C$. This proves the no-new-constants assertion.

Combine the solutions $g_j$ and $U_k$ and invert the unimodular monomial
change. This supplies nonzero solutions to all original diagonal
equations, and they generate the same differential field. Every independent
symbol may be multiplied by an arbitrary nonzero complex constant, giving
$\gm(\C)^d$. In the original coordinates the trivialized integer
combinations must remain fixed, which gives \eqref{eq:torus}. Conversely,
those equations leave exactly the independent $d$ constant multipliers.
\end{proof}

\begin{example}[One phase, two phase generators]
For $a_1=i$ and $a_2=2i$, the obstructions are $\omega$ and $2\omega$.
The lattice is generated by $(2,-1)$, and one can choose solutions with
$y_2=y_1^2$. The Galois group is one-dimensional:
\[
 G=\{(c,c^2):c\in\C^{\times}\}.
\]
\end{example}

\begin{example}[Real dependence does not force an algebraic relation]
For $a_1=i$ and $a_2=\sqrt2\,i$, the obstructions are $\omega$ and
$\sqrt2\,\omega$. They are dependent over $\R$, but have no nonzero
integer relation. Hence the solution extension has two algebraically
independent generators and Galois group $\gm(\C)^2$.
The field exponential on real arguments does not supply a canonical
irrational power of a new oscillatory symbol inside the rational
Picard--Vessiot extension.
\end{example}

\begin{example}[Two logarithmic levels]
For
\[
 a_1=\frac{i}{\omega},\qquad
 a_2=\frac{i}{\omega\log\omega},
\]
the primitive obstructions are $\log\omega$ and $\log\log\omega$.
The first dominates the second by an infinite factor, so no nonzero
integer relation is possible. Again the torus has dimension two. The
smallness of both coefficients does not reduce the number of independent
missing phases.
\end{example}

The theorem treats diagonal systems completely, and supplies useful
information for triangular extensions. It is not a classification of
Picard--Vessiot groups for arbitrary nontriangular matrices over $\E$.

\section{Why the analytic reports do not contradict these results}\label{sec:external}

\subsection{Two commuting derivations before evaluation}
Consider formal power series in a new indeterminate $z$, with coefficients
in $\E$. There are two derivations on this formal series class:
\[
 D_z\left(\sum_{n\geq0}a_nz^n\right)
   =\sum_{n\geq1}na_nz^{n-1},\qquad
 \D_{\mathrm{coeff}}\left(\sum_{n\geq0}a_nz^n\right)
   =\sum_{n\geq0}(\D a_n)z^n.
\]
They commute, and their sum is a derivation. The verification is
coefficientwise and uses only finite convolution sums at each degree.
No evaluation of the formal series is asserted by this construction.

\begin{proposition}[A valid evaluation chain rule]\label[proposition]{prop:eval}
For a polynomial $P(z)\in\E[z]$ and $b\in\E$,
\begin{equation}\label{eq:polynomial-chain}
 \D(P(b))=(\D_{\mathrm{coeff}}P)(b)+(D_zP)(b)\D b.
\end{equation}
In particular, evaluation at $b=\omega$ intertwines the internal
derivation with $\D_{\mathrm{coeff}}+D_z$ on polynomials.
\end{proposition}
\begin{proof}
Expand $P=\sum_{n=0}^{N}a_nz^n$ and apply the finite product rule to
$a_nb^n$. Summing gives the formula.
\end{proof}
If an infinite series is substituted, one must additionally certify its
summability and the derivative manipulations. The polynomial identity does
not make arbitrary formal substitutions valid.

\subsection{The formal exponential at an infinite argument}
The formal series
\[
 F(z)=\sum_{n\geq0}\frac{(iz)^n}{n!}\in\C[[z]]
\]
satisfies $D_zF=iF$. This is an ordinary formal differential equation and
has no relation to the absence of a nonzero element solving $\D y=iy$
until one attempts to evaluate $z$.

At $z=\omega$, the terms have monomials $\omega^n=t^{-n}$. Their additive
valuation support $\{0,-1,-2,\ldots\}$ is not well ordered. Hence this
family has no Hahn sum and does not define an element of $\E$ by the
stated evaluation rule. At an infinitesimal argument the corresponding
series is summable, but the chain rule then involves the derivative of
that infinitesimal argument, not the value $1$.

One may also define global fine-analytic functions by other phase
conventions; the repository and Ehrlich--Kaplan discuss such constructions
\cite{RepoAnalysis,RepoTrig,EK}. Their function-variable derivative does
not establish an internal scalar differential identity after arbitrary
substitution.

\begin{corollary}[A phase-independent exponential obstruction]\label[corollary]{cor:globalexp}
There is no map $\Psi:\E\to\E^{\times}$ satisfying
\begin{equation}\label{eq:globalexpcompat}
 \D\Psi(z)=\Psi(z)\D z\qquad\text{for every }z\in\E.
\end{equation}
No homomorphism, analyticity, or conjugation assumption is needed for this
nonexistence statement.
\end{corollary}
\begin{proof}
At $z=i\omega$, the value $y=\Psi(i\omega)$ would be nonzero and would
satisfy $\D y=iy$, contradicting \cref{cor:noosc}.
\end{proof}
This is stronger in its phase-independence than an obstruction for one
chosen canonical exponential, but is compatible with that existing
repository theorem. It does not deny the existence of a global surcomplex
exponential as a map with other stated properties. It denies the displayed
compatibility with this fixed internal derivation.

\subsection{Germs and composition require their own hypotheses}
The literature on transseries interpreted as surreal function germs
constructs composition on specified classes and also identifies limits to
extending compatible composition to all surreals \cite{BMGerms}. This
context reinforces the need to keep scalar differential structure,
function evaluation, and composition as separate operations. The result
above uses only the simple oscillator obstruction; it does not assume
or reprove a general theorem about all possible notions of composition.

Likewise, the elementary embedding of the transseries field into the
surreals is a theorem about a specified ordered differential language
\cite{ADH}. It should not be silently upgraded to a statement about every
simultaneously added analytic, compositional, or complex exponential
operation. The present article stays within its declared differential
contract and explicit infinitesimal functional calculus.

\section{Set-sized workspaces and an implementable boundary}\label{sec:localization}

\subsection{Differential localization is stronger than algebraic localization}
A Hahn workspace that contains a given collection of coefficients need
not be closed under their internal derivatives, primitives, or phase
operations. The repository's foundational localization perspective
\cite{RepoFoundations} therefore needs a differential refinement for the
present subject.

\begin{proposition}[A set-sized differential workspace]\label[proposition]{prop:localization}
Given a set of surreal and surcomplex input elements, there is a set-sized
real closed subfield $K\subseteq\F$ containing $\R$, $\omega$, and the
real and imaginary parts of all inputs, with the following properties:
\begin{enumerate}[label=(\roman*)]
\item $K$ is closed under $\D$, $\I$, the projections in
\eqref{eq:split}, real exponential, and logarithm on positive elements;
\item $K[i]$ is closed under the infinitesimal exponential and logarithm
whenever their arguments belong to their stated domains;
\item $\D K=K$, the constant fields of $K$ and $K[i]$ are $\R$ and $\C$,
and the bounded-phase and gauge-normal-form theorems hold internally in
$K[i]$, with representatives $K\cap\J$.
\end{enumerate}
\end{proposition}
\begin{proof}
Begin with the set consisting of $\R$, $\omega$, and all real and imaginary
input parts. Form the field it generates. At each subsequent stage adjoin
the values of the following fixed finitary class functions on admissible
tuples from the current stage: $\D$, $\I$, the three normal-form
projections, real exponential and positive logarithm, and the real and
imaginary component functions of $\Ex$ and $\Lo$. Also adjoin all real
roots, in $\F$, of nonzero polynomials with coefficients in the current
field. Each polynomial has finitely many roots. Replacement and union
make each stage a set. Take the field generated by these elements and
repeat for ordinary countably many stages.

The union $K$ is a set-sized field closed under all the specified
operations. Every polynomial over $K$ has its finite coefficient list
inside some stage, so its real roots in $\F$ lie in a later stage. Since
$\F$ is real closed, the resulting ordered subfield is real closed as
well. Thus $K[i]$ is algebraically closed.

Closure under $\I$ gives surjectivity of $\D$ on $K$. Constants remain
exactly $\R$ and $\C$, since they are inherited from the ambient field
and already lie in $K$ and $K[i]$. In the polar decomposition of an element
of $K[i]$, the modulus uses a real root, its logarithm is in $K$, the
ordinary standard part is in $\C$, and the infinitesimal logarithm is
available by construction. All proofs of the image, bounded-phase, and
gauge theorems therefore remain internal to $K[i]$. The projections ensure
that their canonical obstructions lie in $K\cap\J$.
\end{proof}

This argument adjoins values of \emph{specified finitary operations}, some
of which are defined in the ambient field by Hahn sums. It does not assert
closure under all set-indexed summable families. That stronger demand is
unnecessary and would not follow from this construction.

The proposition also allows every finite diagonal or triangular problem
in this article to be treated with an ordinary set-sized differential
field as base. In that setting the usual Picard--Vessiot terminology and
algebraic groups apply without proper-class bookkeeping. The explicit
rational extension and no-new-constants proofs do not require any
unmentioned compactness principle.

\subsection{A practical coefficient-domain contract}
A symbolic system should expose an internal derivative as a separate
operation from the host language's derivative in an independent variable.
For a coefficient $a=p+iq$, the logical workflow is
\begin{equation}\label{eq:workflow}
 a\ \longmapsto\ B=\I q\ \longmapsto\ \up B
 \ \longmapsto\
 \begin{cases}
  \text{a base-field solution},&\up B=0,\\
  \text{a certified obstruction and extension symbol},&\up B\ne0.
 \end{cases}
\end{equation}
This is a mathematical specification, not a universal terminating
algorithm. For arbitrary program-described surreal normal forms,
integration, exact support access, and zero testing may not be available.
The repository's capability distinctions therefore remain necessary
\cite{RepoCAS}.

A first effective tier can use finite Gaussian-rational Laurent expressions
in $t$ and named logarithmic primitives. A larger certified tier can add
real-exponent Hahn coefficients satisfying the explicit support conditions
of \cref{thm:power}. The iterated-logarithm family gives another useful
constructor with a closed-form solvability test. A general backend may
instead return a well-defined symbolic primitive and state that the
purely infinite part has not been decided. It should not report
``no solution'' merely because its representation cannot compute that part.

For rank-one extensions, a formal symbol can carry the defining derivative
$\D U=aU$ together with a certificate that $\ob(a)\ne0$. For several
symbols, the integer relation lattice is additional data: treating
$i\omega$ and $2i\omega$ as independent phases would produce the wrong
Galois group, while treating $i\omega$ and $i\sqrt2\,\omega$ as
algebraically dependent merely because their phases are real multiples
would also be wrong.

\section{Verification, reproducibility, and remaining directions}\label{sec:verification}

\subsection{What was executed}
The accompanying \texttt{code/verify.py} uses only the Python standard
library. Its coefficients are exact Gaussian rationals; its finite
Laurent exponents are rational numbers represented by
\texttt{fractions.Fraction}. The internal derivative is implemented as
$c t^q\mapsto-qc t^{q+1}$. Primitive tests include a separate logarithmic
term at exponent $1$. Infinitesimal exponential identities are checked
as finite jets with a declared cutoff, not by floating-point approximation.

The bundled run used Python 3.13.5, deterministic seed 20260921, and cutoff
16. Its actual recorded result is 1,395 checks passed:
\begin{center}\small
\begin{tabularx}{\textwidth}{@{}Yr@{}}
\toprule
Check category & Count\\
\midrule
Leibniz rule & 250\\
Additivity & 250\\
Primitive identity including the logarithmic term & 250\\
Conjugation compatibility & 250\\
Power-threshold regression cases & 49\\
Positive-phase primitives & 100\\
Small-phase differential equation, as a jet & 100\\
Unit-circle exponential identity, as a jet & 100\\
Cayley logarithmic derivative, as a jet & 1\\
Exact polynomial-forcing inverses & 36\\
Repeated real-root polynomial amplitudes & 9\\
\midrule
Total & 1,395\\
\bottomrule
\end{tabularx}
\end{center}
The script and \texttt{data/verification.json} are supplied so that these
claims can be reproduced. Python version and elapsed time, when reported
by a local environment, need not match the bundled run.

\subsection{What these checks do not establish}
No finite check verifies a proper-class quantifier, Hahn summability for
arbitrary supports, the Berarducci--Mantova construction, a general
no-new-constants theorem, or the novelty of an argument. No Lean development
is included. Mathematical justification for the general statements lies
in the displayed proofs and their explicitly cited published inputs.

The document's repository assessment is revision-specific. It acknowledges
existing results rather than inferring their absence from a missing table
of contents entry. No claim is made that the selected gap remains unfilled
on a later branch or commit. The bibliography is targeted, not an exhaustive
priority search.

\subsection{Useful extensions of this chapter}
The rank-one obstruction and upper-triangular criterion leave several
natural directions. For a general matrix system, the trace condition is
only the first invariant; a useful next step is to understand which
nontriangular differential modules admit a filtration whose rank-one
factors have computable phase obstructions. This would require genuine
module-theoretic information, not merely eigenvalues of an instantaneous
coefficient matrix.

A second direction is effective integration in successively richer
transserial coefficient domains. The power and finite iterated-logarithm
examples show what a certified domain-specific algorithm can deliver.
Extending such algorithms should preserve the distinction between a
representation of a primitive and a decision about its infinite part.

A third direction is formalization. The phase-image theorem needs a
strongly additive derivation, infinitesimal exponential and logarithm,
constant-field identification, and normalized primitive operations. These
can be isolated as interfaces before attempting to formalize the full
surreal derivation. The rational no-new-constants proof can then be checked
as a largely ordinary differential-algebra argument. None of these proposed
formalization steps is claimed to have been carried out here.

\section{Conclusion}
The internal derivative does not turn every surcomplex number into an
ordinary complex-valued function at infinity. Its real exponential
integration is extremely rich, but its unit-circle phases are constrained:
only a constant ordinary phase and an infinitesimal moving phase already
live in the base field. The exact consequence is
\[
 \boxed{\quad\ld(\No[i]^{\times})=\No+i\D\mm\quad}.
\]
The purely infinite part of a primitive of the imaginary coefficient is
therefore the complete rank-one obstruction. It supplies a canonical
normal form, distinguishes algebraic closure from differential solution
closure, explains the missing oscillator, and determines how many formal
phase generators a diagonal Picard--Vessiot extension needs.

This fills a specific conceptual space between the repository's scalar
algebra, finite-angle trigonometry, external-variable analysis, and CAS
capability contracts. Keeping those structures separate is not merely a
notation policy: their compatibility and incompatibility are theorems.

\appendix
\section{Notation and hypotheses at a glance}\label{app:notation}
\begin{longtable}{@{}P{0.19\textwidth}P{0.74\textwidth}@{}}
\toprule
Symbol & Meaning\\
\midrule\endhead
$\F=\No$ & Real surreal field, with the fixed simplest BM derivation.\\[0.35em]
$\E=\F[i]$ & Algebraically closed surcomplex field; constant field $\C$.\\[0.35em]
$\D$ & Internal derivation, $\D\omega=1$; strongly additive, real-exp
compatible, surjective, and small on infinitesimals.\\[0.35em]
$D_z$ & Differentiation in an external variable, holding all surcomplex
coefficients fixed.\\[0.35em]
$\OO,\mm$ & Finite real elements and real infinitesimals.\\[0.35em]
$\J$ & Purely infinite normal forms, including zero; a representative
class, not a class of proper-class cosets.\\[0.35em]
$\up{x},c_0(x),\down{x}$ & Infinite, real constant, and infinitesimal
normal-form projections.\\[0.35em]
$\I b$ & Unique primitive of $b$ with $c_0(\I b)=0$.\\[0.35em]
$\Ex,\Lo$ & Strong Hahn exponential on $\mm_{\C}$ and logarithm on
$1+\mm_{\C}$.\\[0.35em]
$\ld(y)$ & Logarithmic derivative $\D y/y$, for $y\ne0$.\\[0.35em]
$\ob(a)$ & Purely infinite part $\up{(\I\Imm a)}$ of the imaginary
primitive.\\[0.35em]
$t$ & $\omega^{-1}$. The real-exponent identity is
$\D t^\alpha=-\alpha t^{\alpha+1}$, $\alpha\in\R$.\\[0.35em]
$\gm(\C)$ & Multiplicative algebraic group, acting by constant
multiplication on a missing solution symbol.\\
\bottomrule
\end{longtable}

\section{Reproduction instructions}
The source \texttt{article.tex} is standalone: its bibliography is internal,
and no external images, bibliography database, shell escape, or downloaded
font files are required. A standard TeX installation with the packages
listed in the preamble suffices.
\begin{verbatim}
python code/verify.py
latexmk -pdf -interaction=nonstopmode -halt-on-error article.tex
\end{verbatim}
When \texttt{latexmk} is unavailable, run \texttt{pdflatex} three times with
\texttt{-interaction=nonstopmode -halt-on-error}. The supplied shell and
PowerShell scripts automate this fallback. They do not install software or
modify the repository. The archive contains the finished PDF and its source,
not a patch applied to the remote repository.

\clearpage
\phantomsection
\addcontentsline{toc}{section}{References}
\begin{thebibliography}{99}\small\raggedright
\bibitem{BM}
Alessandro Berarducci and Vincenzo Mantova,
\emph{Surreal numbers, derivations and transseries},
Journal of the European Mathematical Society \textbf{20} (2018), 339--390.
\href{https://doi.org/10.4171/JEMS/769}{doi:10.4171/JEMS/769}.
\href{https://arxiv.org/abs/1503.00315}{arXiv:1503.00315}.
Theorems A and B in the introduction, and the underlying Theorems 6.30
and 7.7; normal forms and summability in Sections 2--3.

\bibitem{BMGerms}
Alessandro Berarducci and Vincenzo Mantova,
\emph{Transseries as germs of surreal functions},
Transactions of the American Mathematical Society \textbf{371} (2019),
3549--3592.
\href{https://doi.org/10.1090/tran/7428}{doi:10.1090/tran/7428}.
\href{https://arxiv.org/abs/1703.01995}{arXiv:1703.01995}.
See the composition framework and its limitations, especially Section 8.

\bibitem{ADH}
Matthias Aschenbrenner, Lou van den Dries, and Joris van der Hoeven,
\emph{The surreal numbers as a universal $H$-field},
Journal of the European Mathematical Society \textbf{21} (2019), 1179--1199.
\href{https://doi.org/10.4171/JEMS/858}{doi:10.4171/JEMS/858}.
\href{https://arxiv.org/abs/1512.02267}{arXiv:1512.02267}.
Theorem 1 concerns an elementary embedding in the ordered differential
language.

\bibitem{EK}
Philip Ehrlich and Elliot Kaplan,
\emph{Surreal ordered exponential fields},
Journal of Symbolic Logic \textbf{86}, no. 3 (2021), 1066--1115.
\href{https://doi.org/10.1017/jsl.2021.59}{doi:10.1017/jsl.2021.59}.
\href{https://arxiv.org/abs/2002.07739}{arXiv:2002.07739}.
See Section 11 for global trigonometric and surcomplex exponential
constructions.

\bibitem{PS}
Marius van der Put and Michael F. Singer,
\emph{Galois Theory of Linear Differential Equations},
Grundlehren der mathematischen Wissenschaften, vol. 328, Springer, 2003.
\href{https://doi.org/10.1007/978-3-642-55750-7}{doi:10.1007/978-3-642-55750-7}.
Chapters 1--2 supply the Picard--Vessiot and differential-module terminology.

\bibitem{RepoMap}
Vladimir Reshetnikov, \emph{Surreal repository: documentation map},
\texttt{docs/README.md}, revision
\texttt{e260237db9b71da8b74a0c13c8e6355119091100}, inspected September 2026.
\href{https://github.com/VladimirReshetnikov/Surreal/blob/e260237db9b71da8b74a0c13c8e6355119091100/docs/README.md}{Commit-pinned documentation map}.
The following repository references all use this same snapshot.

\bibitem{RepoAnalysis}
\emph{Surcomplex Analysis: Germs, coherent Hahn sections, and infinitesimal
zero geometry}, repository report,
\path{docs/surcomplex/analysis/article.tex}.
\href{https://github.com/VladimirReshetnikov/Surreal/blob/e260237db9b71da8b74a0c13c8e6355119091100/docs/surcomplex/analysis/article.tex}{Commit-pinned source}.
In particular, the derivative distinction and the theorem with source label
\texttt{b:derivation-obstruction}.

\bibitem{RepoTrig}
\emph{Trigonometry on the Surcomplex Plane: Canonical Phases,
Arbitrary-Scale Geometry, and Infinitesimal Degeneration}, repository report,
\path{docs/surcomplex/trigonometry/article.tex}.
\href{https://github.com/VladimirReshetnikov/Surreal/blob/e260237db9b71da8b74a0c13c8e6355119091100/docs/surcomplex/trigonometry/article.tex}{Commit-pinned source}.
Finite-angle phase decomposition and the distinction from scalar derivations.

\bibitem{RepoCAS}
\emph{Surreal and Surcomplex Numbers in Symbolic Computer Algebra},
repository report,
\path{docs/foundations-and-computation/computer-algebra/article.tex}.
\href{https://github.com/VladimirReshetnikov/Surreal/blob/e260237db9b71da8b74a0c13c8e6355119091100/docs/foundations-and-computation/computer-algebra/article.tex}{Commit-pinned source}.
The section with source label \texttt{cas:sec-derivatives} and the
capability-contract discussion.

\bibitem{RepoGaps}
\emph{A countable genetic cut across the finite--infinite gap},
repository report description,
\path{docs/surreal/genetic-gaps-and-primitives/README.md}.
\href{https://github.com/VladimirReshetnikov/Surreal/blob/e260237db9b71da8b74a0c13c8e6355119091100/docs/surreal/genetic-gaps-and-primitives/README.md}{Commit-pinned description}.
The stated derivative is pointwise order-field differentiation of a function.

\bibitem{RepoFoundations}
\emph{Foundations for surreal and surcomplex mathematics},
repository report description,
\path{docs/foundations-and-computation/foundations/README.md}.
\href{https://github.com/VladimirReshetnikov/Surreal/blob/e260237db9b71da8b74a0c13c8e6355119091100/docs/foundations-and-computation/foundations/README.md}{Commit-pinned description}.
Sets, classes, workspaces, and the limits of verification claims.
\end{thebibliography}
\end{document}
